\documentclass[a4paper,8pt]{extarticle} % 10pt - 30% = 7pt

% Packages
\usepackage[utf8]{inputenc} % Unicode support (Umlauts etc.)
\usepackage{amsmath,mathtools} % Advanced math typesetting
\usepackage{physics} % Physics notation
\usepackage{hyperref} % Add a link to your document
\usepackage{geometry} % Margins
\usepackage{multicol} % for multicolumn layout
\usepackage{xcolor} % For color
\usepackage{tcolorbox} % For colored boxes
\usepackage{fancyhdr} % Headers and footers

\geometry{top=1in, headheight=0.5in, headsep=0.1in, margin=0.5in}

% Compact itemize
\usepackage{enumitem}
\setlist{nosep}

% Compact sections
\usepackage[compact]{titlesec}
\titlespacing{\section}{0pt}{*0}{*0}
\titlespacing{\subsection}{0pt}{*0}{*0}
\titlespacing{\subsubsection}{0pt}{*0}{*0}

% Reduce line spacing
\renewcommand{\baselinestretch}{0.8}

% Color scheme
\definecolor{lightblue}{rgb}{0.75,0.85,1}

% Headers
\pagestyle{fancy}
\fancyhf{} % clear all header and footer fields
\renewcommand{\headrulewidth}{0pt} % no line in header area
\fancyhead[C]{Formelsammlung MMP II} % Center

% New command for creating boxes with black text
\newcommand{\mybox}[2]{
    \begin{tcolorbox}[colback=lightblue!5!white,colframe=lightblue!75!black,boxsep=1pt,arc=0pt,outer arc=0pt,title={\textcolor{black}{#1}}]
        \textcolor{black}{#2}
    \end{tcolorbox}
}

% Start the document
\begin{document}

\fontsize{7pt}{8pt}\selectfont

\begin{multicols}{2}

\mybox{Title}{
Text inside of the box.
With equations:
    \begin{subequations}
    \begin{align}
    v &= u + at \\
    s &= ut + \frac{1}{2}at^2 \\
    v^2 &= u^2 + 2as
    \end{align}
    \end{subequations}
}

\mybox{Dynamics}{
    \begin{equation}
    F = ma
    \end{equation}
}

\mybox{Energy}{
    \begin{equation}
    E = mc^2
    \end{equation}
}

\mybox{Quantum Mechanics}{
    \begin{equation}
    i\hbar\frac{\partial}{\partial t}\Psi = \hat{H}\Psi
    \end{equation}
}

\end{multicols}

% End the document
\end{document}
